\documentclass[11pt]{article}
\usepackage[margin=1in]{geometry}
\usepackage{amsmath,amssymb,mathtools,booktabs,mathrsfs}
\usepackage[hidelinks]{hyperref}
\setlength{\emergencystretch}{3em}
\newcommand{\CC}{\mathcal C}
\newcommand{\RR}{\mathfrak R_T}
\newcommand{\BF}{\mathcal B}
\newcommand{\ZZ}{\mathscr Z}
\newcommand{\norm}[1]{\lVert#1\rVert}
\title{Numerical Companion to the Hardy-Gauge Contour Matrix}
\author{Chaoyu Hu and Jizheng Chen}
\date{25 September 2026}

\begin{document}
\maketitle

\begin{abstract}
This independent numerical note specifies matrices that can be computed before
any stationary-phase or scalar-kernel reduction. It distinguishes the original
closed-contour matrix $A_F$ from the core arithmetic matrix $A_T^{\rm true}$,
gives their exact four-block relation, describes the quadrature and FFT
approximations actually used, and records finite-scale spectra and consistency
errors. The data expose which analytical estimates would be needed in a proof;
they do not establish asymptotic positivity or the density-one hypothesis.
\end{abstract}

\section{The matrix to be approximated}

Fix $T$, choose boundary heights $T_1\in[T,T+1]$ and
$T_2\in[2T,2T+1]$, and let $L=\log(T/2\pi)$. The symmetric rectangle has
$-1\le\Re s\le2$ and $T_1\le\Im s\le T_2$, with counterclockwise boundary.
Use the Hardy gauge $\ZZ=q\zeta$, where $q^2=\chi(1-s)$, and put
\begin{equation}
\begin{split}
f&=-\tfrac12\chi'/\chi,\qquad
P=-\zeta'/\zeta,\qquad Q=(\zeta'/\zeta)',\\
L_1&=Z_1'/Z_1,\qquad Z_1=\zeta'+f\zeta,\qquad H=f+L_1,\\
\CC&=\zeta^2Q,\qquad
C_{\ZZ}^{\rm norm}=\frac{q^2}{L^2}(\CC+f'\zeta^2).
\end{split}
\label{eq:fields}
\end{equation}
For packet centres $\tau_j$ and packet polarization
$\BF_{j k}(s)=B_k(z(s))B_j^\#(z(s))$, $z(s)=-i(s-1/2)$, the original matrix is
\begin{equation}
(A_F)_{jk}=\frac{1}{2\pi i}\oint_{\partial\mathcal R_T^{\rm sym}}
L_1(s)C_{\ZZ}^{\rm norm}(s)\BF_{jk}(s)\,ds.
\label{eq:af}
\end{equation}
The same closed integral with $H$ in place of $L_1$ is equal to
\eqref{eq:af}: the difference is the integral of the holomorphic
$fC_{\ZZ}^{\rm norm}\BF_{jk}$. This equality is a useful independent
numerical check, but it does \emph{not} remove the right-edge Hermitian
$f$ block.

Define the upward right-edge integral $I_D$ by
\begin{equation}
(I_D)_{jk}=\frac{1}{2\pi L^2}\int_{T_1}^{T_2}
q(2+it)^2D(2+it)\BF_{jk}(2+it)\,dt,
\qquad \RR[D;1,1]=I_D+I_D^*.
\label{eq:right}
\end{equation}
The factor $L^{-2}$ is mandatory: without it the resulting matrix has a
different scale from $A_F$. Leg multipliers in $\RR[D;a,b]$ act on the two
packet syntheses before integration, as in the main manuscript.

\section{The exact four blocks}

On the right edge, set
\begin{equation}
H_{\rm ar}=-P+\frac{Q}{f-P},\qquad
D_{\rm ar}=\CC H_{\rm ar},\qquad
D_f=f(\CC+f'\zeta^2),
\label{eq:ar}
\end{equation}
and
\begin{equation}
D_{f'}=\frac{f'}{f-P}\CC+f'\zeta^2H_{\rm ar}
          +\frac{(f')^2}{f-P}\zeta^2.
\label{eq:fp}
\end{equation}
The identity $H=f+H_{\rm ar}+f'/(f-P)$ gives, pointwise on this edge,
\begin{equation}
HC_{\ZZ}^{\rm norm}
=\frac{q^2}{L^2}(D_{\rm ar}+D_f+D_{f'}).
\label{eq:pointwise}
\end{equation}
Let $\vartheta_c$ be the chosen core multiplier. The four Hermitian blocks
are defined on the \emph{same full packet space} by
\begin{align}
A_T^{\rm true}&=\RR[D_{\rm ar};\vartheta_c,\vartheta_c],
\label{eq:true}\\
L_T^{\rm geom}&=\RR[D_{\rm ar};1,1]-A_T^{\rm true},
\label{eq:geom}\\
J_T^{\rm arch}&=\RR[D_f;1,1],
\label{eq:arch}\\
R_T^{\rm exact}&=\RR[D_{f'};1,1]+R_T^{\rm hor,H},
\label{eq:remainder}
\end{align}
where $R_T^{\rm hor,H}$ is the oriented integral of
$HC_{\ZZ}^{\rm norm}\BF$ on the top and bottom sides. Reflection converts the
left $H$ integral into the adjoint of the right $H$ integral. Consequently,
\begin{equation}
\boxed{A_F=A_T^{\rm true}+L_T^{\rm geom}
                +J_T^{\rm arch}+R_T^{\rm exact}.}
\label{eq:four}
\end{equation}
The equality is algebraic and survives common compression to a retained
subspace. The $ff'\zeta^2$ term belongs to $J_T^{\rm arch}$; omitting it would
make \eqref{eq:four} inexact. The $f'$ source and horizontal contour are grouped
into one fourth block only for bookkeeping. No smallness or sign is assigned to
any block by \eqref{eq:four}.

\section{Finite packet model and discretization}

The numerical specialization uses an exact lattice
$\tau_j=2\pi j/L_{\rm pkt}$ within $[T,2T]$, where
\begin{equation}
L_{\rm pkt}=L+0.1\log L,\qquad
W_T=L-0.1\log L,\qquad u_0=-L/2.
\label{eq:packetparameters}
\end{equation}
The full packet window equals $L_{\rm pkt}^{-1/2}$ on
$[-L_{\rm pkt}/2,L_{\rm pkt}/2]$ and has $C^\infty$ caps of width $0.1$;
the sharp core is $[u_0,u_0+W_T]$. These are explicit finite test parameters,
not the manuscript's eventual asymptotic retained projection.

On a vertical side $s=\sigma+it$, the flat part of a packet leg has the
closed form
\begin{equation}
\begin{aligned}
a_\pm&=\pm(\sigma-1/2)+i(t-\tau_j),\\
P_\pm(t,\tau_j;\sigma,u,w)
&=\frac1{\sqrt{L_{\rm pkt}}}\int_u^{u+w}e^{a_\pm x}\,dx\\
&=\frac{w}{\sqrt{L_{\rm pkt}}}
e^{a_\pm(u+w/2)}\frac{\sinh(a_\pm w/2)}{a_\pm w/2}.
\end{aligned}
\label{eq:packetclosed}
\end{equation}
The quotient is evaluated by its Taylor series near zero. Smooth caps are
integrated with 64-point Gauss--Legendre quadrature. The matrix assembly has
the weighted-product form
\begin{equation}
I_D\simeq P_-^*\operatorname{diag}(w_rF_D(t_r))P_+,
\qquad F_D(t)=\frac{q(2+it)^2D(2+it)}{2\pi L^2},
\label{eq:product}
\end{equation}
with trapezoidal weights $w_r$ on vertical sides. The top and bottom use
Gauss--Legendre nodes in $\sigma$ and their true boundary orientations.
For the direct $A_F$ computation, $\zeta,\zeta',\zeta''$ are evaluated using
SciPy's complex zeta function and fourth-order central differences; $f'$ is
differenced similarly. Thus the calculations here remain floating-point
approximations, including their special-function evaluations.

For the \emph{core right-edge} matrix, the $u$ integral in
\eqref{eq:packetclosed} can instead be sampled and transformed by FFT.
With $t_r-\tau_j$ on a grid of spacing $h_t$, choose $N_{\rm FFT}$ and
$h_u=2\pi/(N_{\rm FFT}h_t)$. A midpoint sum in $u$ gives all packet
frequencies from one inverse FFT, with the origin phase restored explicitly.
The final partial cell is clipped to the exact core endpoint. The direct
closed-form packet integral checks the FFT approximation separately from the
height quadrature. The current \emph{full-contour} script uses the analytic
packet formula and Gauss quadrature, not an FFT for its horizontal sides.

\section{Finite spectra and decomposition checks}

The following rows use the same smooth packet and good-height contour within
each value of $T$. Let $S_0=A_T^{\rm true}$,
$S_1=S_0+L_T^{\rm geom}$, $S_2=S_1+J_T^{\rm arch}$, and
$S_3=S_2+R_T^{\rm exact}$. The direct column comes from the original
$L_1$ closed integral \eqref{eq:af}, independently of the split.

\begin{center}
\begin{tabular}{@{}r l r r@{}}
\toprule
$T$ & matrix & negative eigenvalues & smallest eigenvalue \\
\midrule
80 & $S_0$ & 19 & $-3.303704$ \\
   & $S_1$ & 18 & $-5.292936$ \\
   & $S_2$ & 0  & $\phantom{-}0.597878$ \\
   & $S_3$ & 0  & $\phantom{-}0.592901266$ \\
   & direct $A_F$ & 0 & $\phantom{-}0.592901060$ \\
\midrule
160 & $S_0$ & 46 & $-6.320577$ \\
    & $S_1$ & 46 & $-8.974877$ \\
    & $S_2$ & 2  & $-11.558592$ \\
    & $S_3$ & 0  & $\phantom{-}0.003758395$ \\
    & direct $A_F$ & 0 & $\phantom{-}0.003757118$ \\
\bottomrule
\end{tabular}
\end{center}

The code selects these heights by maximizing the smallest sampled $|Z_1|$
along each horizontal side; this sampling is not a proof that a side is
zero-free. At $T=80$ the contour is $[80.1,160.95]$ and the dimension is 34. The
vertical mesh is $0.0075$ and each horizontal side has 2048 Gauss nodes.
At $T=160$ the contour is $[160.95,320.3]$ and the dimension is 85; the
vertical mesh is $0.015$ and each horizontal side has 1024 nodes. The
archimedean block has Frobenius norms $33.89237$ and $67.86938$, respectively.
It is positive in the first sample but itself has two negative eigenvalues
in the second. The fourth block is essential at $T=160$: adding $J_T^{\rm arch}$
alone does not remove the last two negative directions.

The relative Frobenius discrepancy
$\norm{S_3-A_F^{(H)}}_F/\norm{A_F^{(H)}}_F$ is
$3.34\cdot10^{-7}$ at $T=80$ and $4.82\cdot10^{-6}$ at $T=160$.
Here $A_F^{(H)}$ is the independently assembled closed integral with $H$.
The independent $L_1$ and $H$ closed integrals differ relatively by
$5.43\cdot10^{-6}$ and $1.27\cdot10^{-4}$, respectively. These discrepancies
come from quadrature, finite differences, and floating-point evaluation; the
symbolic identity \eqref{eq:four} has no such error.

The mesh sequence below holds the packet parameters and selected contour
fixed at each $T$. Here $\delta_{\rm split}$ is the relative Frobenius
discrepancy between $S_3$ and the direct $H$ integral, while
$\delta_{L_1,H}$ compares the independently discretized closed integrals.
\begin{center}
\small
\begin{tabular}{@{}r r r r r r@{}}
\toprule
$T$ & vertical step & horizontal nodes & $\lambda_{\min}(A_F)$
& $\delta_{\rm split}$ & $\delta_{L_1,H}$ \\
\midrule
80  & $0.0300$ & 512  & $0.592889$ & $7.09\cdot10^{-7}$ & $8.76\cdot10^{-5}$ \\
80  & $0.0150$ & 1024 & $0.592899$ & $5.35\cdot10^{-7}$ & $2.19\cdot10^{-5}$ \\
80  & $0.0075$ & 2048 & $0.592901$ & $3.34\cdot10^{-7}$ & $5.43\cdot10^{-6}$ \\
160 & $0.0300$ & 512  & $0.006842$ & $6.77\cdot10^{-6}$ & $5.06\cdot10^{-4}$ \\
160 & $0.0150$ & 1024 & $0.003757$ & $4.82\cdot10^{-6}$ & $1.27\cdot10^{-4}$ \\
\bottomrule
\end{tabular}
\end{center}
The $T=160$ minimum moves appreciably on refinement, so its positive sign is
not numerically certified. In the separate core FFT experiment on $[T,2T]$,
the full FFT matrix differs from a matrix using analytic packet integrals by
$1.32\cdot10^{-5}$ in relative Frobenius norm at $T=80$; that comparison
isolates the packet-transform approximation and uses different boundary
heights from the closed-contour table.

For computed Hermitian matrices $A$ and $B$, Weyl's inequality gives
\begin{equation}
|\lambda_{\min}(A)-\lambda_{\min}(B)|
\le\norm{A-B}_2\le\norm{A-B}_F.
\label{eq:weyl}
\end{equation}
This is a check between discretized matrices, not an error bound between a
discretization and the exact contour integral. At $T=160$ the observed
minimum is only about $0.0038$, so its sign is sensitive to a much tighter
quadrature and special-function error analysis. Neither sample proves that
$A_F\succeq0$ at other heights.

\section{What the numerical study should test next}

The paper's conditional hypothesis is about a \emph{whitened one-sided defect},
not the raw negative index of $A_T^{\rm true}$. On the precisely specified
retained space and with its reference Gram $G_T$, the quantity to measure is
\begin{equation}
\beta_-(T;\lambda_T)=
\frac{\operatorname{tr}\!\left[
\left(G_T^{-1/2}(A_T^{\rm true}-\lambda_TP_T)G_T^{-1/2}\right)_-^2
\right]}{(\kappa_*\lambda_T)^2N_T},
\qquad\kappa_*=0.02.
\label{eq:beta}
\end{equation}
It cannot be inferred from the sign count in the table. A finite comparison
on the mean-zero space alone has been computed, but the full retained
projection and a stable treatment of small eigenvalues of $G_T$ are still
required before extrapolating it.

The immediate numerical tasks are to (i) refine vertical, horizontal,
derivative, and special-function precision independently; (ii) repeat the
four-block calculation across contours and packet parameters; (iii) calculate
the spectrum of $L_T^{\rm geom}+J_T^{\rm arch}$ after the exact geometric edge
extraction; and (iv) compute \eqref{eq:beta} after implementing the paper's
complete retained projection. A rigorous sign certificate would also need
explicit enclosures for the quadrature and field-evaluation errors. The
four-block sum is the object to validate before studying any proposed
asymptotic simplification.

\section*{Reproduction}

The tracked scripts are \path{scripts/compute_original_contour.py},
\path{scripts/compute_true_matrix.py}, and
\path{scripts/compare_true_reference.py}. For the first row of the
table, run from the repository root:
\begin{verbatim}
python3 scripts/compute_original_contour.py --height 80 \
  --smooth-cap 0.1 --height-step 0.0075 \
  --horizontal-nodes 2048 --derivative-step 0.001 \
  --output output/numerics/original_T80_smooth_superfine.json \
  --matrix-output output/numerics/original_T80_smooth_superfine.npz
\end{verbatim}
For $T=160$ use \texttt{--height 160 --height-step 0.015
--horizontal-nodes 1024}, with distinct output names. The core FFT is checked
against analytic packet integrals by
\begin{verbatim}
python3 scripts/compute_true_matrix.py --height 80 --a0 0.1 \
  --packet-reserve 0.1 --height-oversample 24 --u-step 0.001 \
  --derivative-step 0.001 --direct-check
\end{verbatim}
The directed gates are \texttt{npm test -- original-contour true-matrix}.
The JSON/NPZ output is locally generated and is not part of the manuscript.

\end{document}
